\section{From pinned loops to a magnetized equilibrium state}\label{sec:kms}

Theorem~\ref{thm:pin-bootstrap} says that, under boundary pinning,
each slot has small probability of belonging to a cycle that misses the
boundary.  We now turn that statement into a state for the original
zero-field dynamics.  First we obtain a tail bound for the exact spectral
distribution of total magnetization in a free box.  A field tending to
zero selects that tail, and translation averages produce a state with
positive magnetization.  The last step verifies both boundary values of
the KMS strip for that same state.

Fix \(d\geq3\) and \(S\in\{\tfrac12,1,\tfrac32,\ldots\}\).
Let \(p_*\) be the threshold supplied by Theorem~\ref{thm:pin-bootstrap}
for \(\ell=2S\), and put \(p=p_*\) and
\(\beta_0(d,S)=5p_*^{-11/10}\).  Throughout this Section,
\(\beta\geq\beta_0(d,S)\) is fixed and finite.  Write
\[
 \Lambda_N=[-N,N]^d\cap\Z^d,\qquad
 V_N=|\Lambda_N|,\qquad
 M_N=\sum_{x\in\Lambda_N}S_x^z,\qquad H_N=H_{\Lambda_N}.
\]
The quantity \(V_N\) here is the number of sites; the slot set remains
\(V_{\Lambda_N}=\Lambda_N\times\{1,\ldots,\ell\}\), with \(\ell=2S\).
Let \(\mu_N\) be the spectral distribution of \(M_N\) in the free
zero-field Gibbs state:
\[
 \mu_N(I)=
 \frac{\Tr\bigl(e^{-\beta H_N}\1_I(M_N)\bigr)}
      {\Tr(e^{-\beta H_N})}
 \quad\text{for Borel sets }I\subset\R.
\]
By Proposition~\ref{prop:spin-loops}, this is exactly the time-zero
magnetization distribution in the cycle-coloring representation.

\subsection{A spectral tail with surface-order cost}

\begin{lemma}[Boundary coloring and a spectral tail]\label{lem:surface-cost}
For every \(N\geq1\), let \(G_N\) be the event in the zero-field
mark-and-color law that every cycle meeting \(D_0\) is colored up.  There
is a constant \(C_{\beta,d,S}<\infty\), independent of \(N\), such that
\begin{align}
 \Pr(G_N)&\geq
       \exp\{-C_{\beta,d,S}|\partial\Lambda_N|\},
       \label{eq:boundary-color-cost}\\
 \mu_N([SV_N/2,\infty))&\geq
       \tfrac12\exp\{-C_{\beta,d,S}|\partial\Lambda_N|\}.
       \label{eq:free-spectral-tail}
\end{align}
Conditionally on \(G_N\), the uncolored picture has law
\(\mathbb P_{D_0}\), and the cycles missing \(D_0\) still have
independent fair colors.
\end{lemma}

\begin{proof}
Let \(n_0\) be the number of cycles meeting \(D_0\), and let \(J_N\)
count inter-site marks incident to a boundary site during one period.
Every cycle meeting the boundary either contains a time-zero boundary
slot or uses an inter-site mark incident to the boundary.  Indeed, in
the absence of a time-zero boundary slot on that cycle, a boundary visit
must begin by a spatial jump; within-site endpoint permutations do not
change the site.  Each mark has two strands and can witness at most two
cycles.  Consequently, picture by picture,
\[
 n_0\leq \ell|\partial\Lambda_N|+2J_N.
\]
There are at most \(2d\ell^2|\partial\Lambda_N|\) slot edges incident
to the boundary.  Their base rate is \(1/2\).  The first factorial
bound of Lemma~\ref{lem:factorial}, applied with no pins, therefore gives
\[
 \mathbb E_{\varnothing}J_N
 \leq 2d\beta\ell^2|\partial\Lambda_N|,
 \qquad
 \mathbb E_{\varnothing}n_0
 \leq(\ell+4d\beta\ell^2)|\partial\Lambda_N|.
\]
Given a picture, \(G_N\) has probability \(2^{-n_0}\).  Jensen's
inequality now proves \eqref{eq:boundary-color-cost}, for example with
\(C_{\beta,d,S}=(\log2)(\ell+4d\beta\ell^2)\).

Multiplying the zero-field picture weight \(2^{\#\mathrm{cycles}}\)
by \(2^{-n_0}\) leaves precisely the weight for cycles missing
\(D_0\).  This proves the assertion about the conditional picture;
conditioning fixes only the colors of the other cycles.  For a
time-zero slot \(i\), its conditional mean spin is thus
\[
 \frac12\bigl(1-\mathbb P_{D_0}(i\in\mathcal B_{D_0})\bigr).
\]
For an unpinned slot the probability on the right is at most \(p\) by
the singleton case of Theorem~\ref{thm:pin-bootstrap}; for a boundary
slot it is zero.  Summing over the \(\ell V_N\) slots gives
\[
 \mathbb E\left[\frac{M_N}{V_N}\,\middle|\,G_N\right]
       \geq S(1-p).
\]
Here \(M_N\) denotes the random variable in the coloring
representation.  If
\(q=\Pr(M_N/V_N\geq S/2\mid G_N)\), its upper bound \(M_N/V_N\leq S\)
implies
\[
 S(1-p)\leq Sq+\frac S2(1-q),\qquad q\geq1-2p>\frac12.
\]
Multiply by \(\Pr(G_N)\) and use the exact spectral interpretation
of the unconditioned coloring law to obtain
\eqref{eq:free-spectral-tail}.
\end{proof}

The boundary event has not changed the Hamiltonian whose spectral law
appears in \eqref{eq:free-spectral-tail}.  It supplies a tail event in
the free Gibbs state whose logarithmic cost is of order the boundary
size.  We next select this tail by a homogeneous field whose volume
contribution is larger than that cost, although the field itself tends
to zero.  The relation between the zero-field magnetization
distribution and field-selected spontaneous magnetization goes back to
Griffiths~\cite{Griffiths1966}; see also
\cite[Section~1]{DysonLiebSimon1978}.  Here the surface-order tail
permits the explicit volume-dependent field used below.

\begin{proposition}[An explicit vanishing-field limit]\label{prop:magnetized-limit}
For \(N\geq1\), set
\[
 h_N=N^{-1/2},\qquad K_N=H_N-h_NM_N,
 \qquad
 \rho_N^{\mathrm{in}}(C)=
   \frac{\Tr(e^{-\beta K_N}C)}{\Tr(e^{-\beta K_N})}
 \quad(C\in\mathcal A_{\Lambda_N}).
\]
Extend \(\rho_N^{\mathrm{in}}\) to a state \(\rho_N\) on
\(\mathcal A\) by the product normalized trace outside \(\Lambda_N\),
and define
\begin{equation}\label{eq:translation-average}
 \omega_N(A)=\frac1{V_N}\sum_{x\in\Lambda_N}\rho_N(T_x(A)).
\end{equation}
There is a subsequence \(N_j\to\infty\) on which \(\omega_{N_j}\)
converges pointwise on \(\mathcal A\) to a translation-invariant state
\(\omega\), and every such subsequential limit satisfies
\begin{equation}\label{eq:positive-limit-magnetization}
 \omega(S_0^z)\geq S/4.
\end{equation}
\end{proposition}

\begin{proof}
Each interaction term commutes with the sum of the two endpoint
\(z\)-spins.  This follows directly from
\[
 \mathbf S_x\cdot\mathbf S_y
 =S_x^zS_y^z+\tfrac12(S_x^+S_y^-+S_x^-S_y^+).
\]
Thus \([H_N,M_N]=0\).  For any \(h>0\), the spectral law in the
Gibbs state of \(H_N-hM_N\) is exactly
\[
 \mu_{N,h}(\dd m)
 =\frac{e^{\beta hm}\mu_N(\dd m)}
        {\int e^{\beta hu}\mu_N(\dd u)}.
\]
Bound the numerator on \(m<SV_N/4\), and retain only
\(m\geq SV_N/2\) in the denominator.  Lemma~\ref{lem:surface-cost}
then gives
\begin{equation}\label{eq:field-tail}
 \mu_{N,h}(({-\infty},SV_N/4))
 \leq 2\exp\left\{
       C_{\beta,d,S}|\partial\Lambda_N|-\frac{\beta hS}{4}V_N
       \right\}.
\end{equation}
For these cubes,
\[
 \frac{|\partial\Lambda_N|}{V_N}
 =1-\left(\frac{2N-1}{2N+1}\right)^d
 \leq\frac{2d}{2N+1}.
\]
At \(h=h_N\), the exponent in \eqref{eq:field-tail} tends to
\(-\infty\): its negative term has order \(N^{d-1/2}\), whereas
its positive term has order \(N^{d-1}\).  Hence, writing
\(q_N=\mu_{N,h_N}(({-\infty},SV_N/4))\), we have \(q_N\to0\).
The spectral bounds \(-S\leq M_N/V_N\leq S\) imply
\begin{equation}\label{eq:field-mean}
 \rho_N^{\mathrm{in}}(M_N/V_N)
 \geq\frac S4(1-q_N)-Sq_N
 =\frac S4-\frac{5S}{4}q_N.
\end{equation}
In particular this estimate uses a single explicit volume-field
sequence for the fixed \(\beta\).

Choose a countable norm-dense set in the union of local matrix
algebras, closed under rational complex linear combinations and
adjoints.  The bound \(|\omega_N(A)|\leq\|A\|\) permits diagonal
subsequence selection on that set.  Norm approximation then gives
pointwise convergence on \(\mathcal A\).  The resulting functional
\(\omega\) is linear, normalized and positive, since these properties
hold for each \(\omega_N\) and pass to pointwise limits.  It is
therefore a state.

For every fixed \(y\in\Z^d\) and \(A\in\mathcal A\),
\[
 |\omega_N(T_y(A))-\omega_N(A)|
 \leq\frac{|(\Lambda_N+y)\mathbin\triangle\Lambda_N|}{V_N}\|A\|
 \longrightarrow0.
\]
This proves translation invariance of \(\omega\).  Finally,
\[
 \omega_N(S_0^z)
 =\frac1{V_N}\sum_{x\in\Lambda_N}\rho_N(S_x^z)
 =\rho_N^{\mathrm{in}}(M_N/V_N).
\]
Taking the limit in \eqref{eq:field-mean} proves
\eqref{eq:positive-limit-magnetization}.
\end{proof}

\subsection{Dynamics far from the translated boundary}\label{subsec:deep-dynamics}

Fix from now on one state subsequence \(N_j\) provided by
Proposition~\ref{prop:magnetized-limit}.  We will prove the KMS
condition along this same subsequence for every local pair; no further
subsequence depending on that pair will be taken.  The needed
dynamical approximation must hold uniformly over translations whose
distance from the box boundary diverges.

We first check the hypotheses of the dynamics existence theorem cited in
Section~\ref{sec:model}, \cite[Theorem~2.2]{NachtergaeleOgataSims2006}.
Use the lattice graph metric and the summable function
$F(r)=(1+r)^{-d-1}$, which has the convolution bound required there.
The interaction is supported on nearest-neighbor pairs, each term has
norm at most $3S^2$, and every site belongs to $2d$ such pairs.
Thus its interaction norm in that theorem is finite, with parameter
$a=0$.

For a finite cube \(\Gamma\), write
\(\tau_t^\Gamma(B)=e^{itH_\Gamma}Be^{-itH_\Gamma}\), acting trivially
outside \(\Gamma\).  Likewise let
\[
 \alpha_z^N(B)=e^{izK_N}Be^{-izK_N}\qquad(z\in\C).
\]
For real \(t\), these maps are isometric automorphisms.  For complex
\(z\), \(\alpha_z^N\) is defined on the entire quasi-local algebra
because \(K_N\) is a bounded local operator.

\begin{lemma}[Uniform approximation in deep translated boxes]\label{lem:deep-dynamics}
Let \(B\) be a local observable, let \(L\) be a finite set containing
\(\supp B\) and the origin, and let \(0<T<\infty\).  Then
\begin{equation}\label{eq:deep-box-zero-field}
 \lim_{r\to\infty}\ 
 \sup_{\substack{\Gamma\text{ a finite cube}\\
                 \dist_\infty(L,\Gamma^c)\geq r}}
 \ \sup_{|t|\leq T}\|\tau_t^\Gamma(B)-\tau_t(B)\|=0.
\end{equation}
Here \(\Gamma^c=\Z^d\setminus\Gamma\).  If
\[
 I_N=\{x\in\Lambda_N:
       \dist_\infty(x,\Lambda_N^c)\geq\lfloor\sqrt N\rfloor\},
\]
then \(|\Lambda_N\setminus I_N|/V_N\to0\) and
\begin{equation}\label{eq:deep-box-with-field}
 \lim_{N\to\infty}\sup_{x\in I_N}\sup_{|t|\leq T}
 \|\alpha_t^N(T_x(B))-T_x(\tau_t(B))\|=0.
\end{equation}
\end{lemma}

\begin{proof}
First fix \(t\).  If \eqref{eq:deep-box-zero-field} failed at that
time, there would be \(\varepsilon>0\) and cubes \(\Gamma_k\)
whose distances from \(L\) to their complements tend to infinity,
with
\(\|\tau_t^{\Gamma_k}(B)-\tau_t(B)\|\geq\varepsilon\).
Such cubes contain centered boxes with radii tending to infinity.
We can therefore select a subsequence that is increasing by inclusion
and exhausts \(\Z^d\): after any finite selected cube, a sufficiently
late cube contains it and the next prescribed centered box.  This
contradicts the norm convergence along increasing boxes in the
definition of \(\tau_t\).

To make this convergence uniform in time, write
\(\Phi(\{u,v\})=-\mathbf S_u\cdot\mathbf S_v\).  Only edges meeting
\(L\) contribute to \([H_\Gamma,B]\), so for every \(\Gamma\)
\[
 \|[H_\Gamma,B]\|\leq
 2\|B\|\sum_{\substack{e\text{ nearest-neighbor edge}\\e\cap L\ne\varnothing}}
                   \|\Phi(e)\|=:c_B<\infty.
\]
Consequently
\(\|\tau_t^\Gamma(B)-\tau_u^\Gamma(B)\|\leq c_B|t-u|\).
Passing to the infinite-volume limit gives the same bound for
\(\tau_t(B)\).  A finite mesh of \([-T,T]\), together with the
fixed-time conclusion, proves \eqref{eq:deep-box-zero-field}.

For \(x\in I_N\), set \(\Gamma=\Lambda_N-x\).  Then
\(\dist_\infty(L,\Gamma^c)\to\infty\) uniformly over these
\(x\), since \(L\) is fixed.  For all sufficiently large \(N\),
\(L\subset\Gamma\).  Define the local rotation
\[
 \gamma_\theta^L(B)
   =e^{i\theta\sum_{u\in L}S_u^z}B
      e^{-i\theta\sum_{u\in L}S_u^z}.
\]
The commutation \([H_N,M_N]=0\) and translation of the finite
Hamiltonian give the exact formula
\begin{equation}\label{eq:field-rotation-sign}
 T_{-x}\bigl(\alpha_t^N(T_x(B))\bigr)
       =\tau_t^\Gamma\bigl(\gamma_{-h_Nt}^L(B)\bigr).
\end{equation}
The norm effect of this rotation is bounded by
\[
 \|\gamma_{-h_Nt}^L(B)-B\|
 \leq h_N|t|\left\|\left[\sum_{u\in L}S_u^z,B\right]\right\|.
\]
Use the isometry of \(\tau_t^\Gamma\), followed by
\eqref{eq:deep-box-zero-field}, to prove
\eqref{eq:deep-box-with-field}.  Finally, the excluded sites lie in
a layer of thickness at most \(\lfloor\sqrt N\rfloor\), whose
fraction of the cube is \(O(N^{-1/2})\).
\end{proof}

\subsection{The analytic strip and its two boundary values}

We will use the following elementary strip lemma.  Its first part
establishes a common strip bound from finite, possibly nonuniform
interior bounds.  Only after obtaining that common bound does its
second part use Gaussian decay to control convergence on the whole
boundary.

\begin{lemma}[Bounded strip limits]\label{lem:strip-limit}
Fix \(\beta>0\), and put
\(\mathcal S_\beta=\{z\in\C:0\leq\operatorname{Im}z\leq\beta\}\).
Let \((F_n)_{n\geq1}\) be functions continuous on
\(\mathcal S_\beta\), analytic in its interior, and individually
bounded there.  Suppose a constant \(C\geq0\) satisfies
\[
 |F_n(t)|\leq C,\qquad |F_n(t+i\beta)|\leq C
 \quad(n\geq1,\ t\in\R).
\]
Then \(|F_n(z)|\leq C\) throughout \(\mathcal S_\beta\).
If the two boundary sequences converge locally uniformly on \(\R\)
to \(f_0\) and \(f_\beta\), respectively, then the full sequence
\(F_n\) converges locally uniformly on the closed strip to a bounded
function \(F\), continuous on the closed strip and analytic in its
interior, with
\[
 |F(z)|\leq C,\qquad F(t)=f_0(t),\qquad F(t+i\beta)=f_\beta(t).
\]
\end{lemma}

\begin{proof}
Fix \(n\), and let \(M_n<\infty\) bound \(|F_n|\) on the strip.
For \(\varepsilon>0\), apply the maximum principle to
\(e^{-\varepsilon z^2}F_n(z)\) on
\([-R,R]+i[0,\beta]\).  Its modulus is at most
\(Ce^{\varepsilon\beta^2}\) on the horizontal sides and at most
\(M_ne^{-\varepsilon R^2+\varepsilon\beta^2}\) on the vertical
sides.  Send \(R\to\infty\) while holding a point \(z\) fixed,
then send \(\varepsilon\downarrow0\).  This proves
\(|F_n(z)|\leq C\).  Only the already assumed finite bound
\(M_n\) was needed to remove the vertical sides.

Now set \(G_n(z)=e^{-z^2}F_n(z)\).  The common bound just proved
implies
\[
 |G_n(t+iy)|\leq Ce^{-t^2+y^2}
       \quad(0\leq y\leq\beta).
\]
Local uniform convergence of the two boundary sequences and the
Gaussian tail therefore make \((G_n)\) uniformly Cauchy on both
horizontal edges.  Write
\[
 \delta_{n,m}=
 \max_{y\in\{0,\beta\}}\sup_{t\in\R}
           |G_n(t+iy)-G_m(t+iy)|;
 \qquad \delta_{n,m}\longrightarrow0.
\]
The maximum principle on the same rectangles, now applied to
\(G_n-G_m\), has horizontal bound \(\delta_{n,m}\) and vertical
bound \(2Ce^{-R^2+\beta^2}\).  Letting \(R\to\infty\) shows
\[
 \sup_{z\in\mathcal S_\beta}|G_n(z)-G_m(z)|\leq\delta_{n,m}.
\]
Hence \(G_n\) converges uniformly on the closed strip to a
continuous function \(G\), analytic in the interior.  Multiplication
by \(e^{z^2}\) gives local uniform convergence of \(F_n\) on the
closed strip to \(F=e^{z^2}G\).  Its claimed boundary values and
the bound \(|F|\leq C\) follow by taking pointwise limits.
\end{proof}

\begin{proof}[Completion of the proof of Theorem~\ref{thm:main}]
Let \(\omega\) and \(N_j\) be as in
Proposition~\ref{prop:magnetized-limit}.  For local observables
\(A,B\), define the entire functions
\begin{equation}\label{eq:finite-strip-functions}
 F_N^{A,B}(z)=\frac1{V_N}\sum_{x\in\Lambda_N}
       \rho_N\bigl(T_x(A)\alpha_z^N(T_x(B))\bigr).
\end{equation}
Trace cyclicity gives both identities
\begin{align}
 F_N^{A,B}(t)
 &=\frac1{V_N}\sum_{x\in\Lambda_N}
       \rho_N\bigl(T_x(A)\alpha_t^N(T_x(B))\bigr),
       \label{eq:finite-bottom-edge}\\
 F_N^{A,B}(t+i\beta)
 &=\frac1{V_N}\sum_{x\in\Lambda_N}
       \rho_N\bigl(\alpha_t^N(T_x(B))T_x(A)\bigr).
       \label{eq:finite-top-edge}
\end{align}
These formulas also hold when the translated supports cross the
boundary of \(\Lambda_N\).  To see this without a support
restriction, take a finite set containing \(\Lambda_N\) and both
translated supports.  On its matrix algebra the density of \(\rho_N\)
is \(e^{-\beta K_N}/\Tr(e^{-\beta K_N})\) tensored with a normalized
identity.  Equivalently, it is the Gibbs density for the operator
\(K_N\otimes I\) on that larger algebra.  The dynamics is generated
by that same operator, so finite-dimensional trace cyclicity proves
\eqref{eq:finite-top-edge} there as well.

For real \(t\), the isometry of \(\alpha_t^N\) bounds both
right-hand sides by \(\|A\|\|B\|\).  Before using any common strip
bound, we also have the finite bound
\[
 |F_N^{A,B}(t+iy)|
 \leq \|A\|\|B\|e^{2\beta\|K_N\|}
 \quad(t\in\R,\ 0\leq y\leq\beta).
\]
Indeed, the real-time unitary factors have norm one, and
\(\|e^{\pm yK_N}\|\leq e^{\beta\|K_N\|}\).
The first part of Lemma~\ref{lem:strip-limit} therefore applies and
gives the uniform bound
\begin{equation}\label{eq:strip-product-bound}
 |F_N^{A,B}(z)|\leq\|A\|\|B\|
       \quad(z\in\mathcal S_\beta,\ N\geq1).
\end{equation}

We next identify the two boundary limits.  Given \(0<T<\infty\),
Lemma~\ref{lem:deep-dynamics} supplies numbers \(\varepsilon_N(T)\to0\)
bounding the norm error in \eqref{eq:deep-box-with-field} on \(I_N\).
On its complement the error is at most \(2\|B\|\), since both
real-time dynamics are isometric.  It follows that, uniformly for
\(|t|\leq T\),
\begin{align}
 \left|F_N^{A,B}(t)-\omega_N(A\tau_t(B))\right|
 &\leq\|A\|\varepsilon_N(T)
   +2\|A\|\|B\|\frac{|\Lambda_N\setminus I_N|}{V_N},
       \label{eq:bottom-edge-error}\\
 \left|F_N^{A,B}(t+i\beta)-\omega_N(\tau_t(B)A)\right|
 &\leq\|A\|\varepsilon_N(T)
   +2\|A\|\|B\|\frac{|\Lambda_N\setminus I_N|}{V_N}.
       \label{eq:top-edge-error}
\end{align}
For example, the comparison term in the first line is exactly
\(V_N^{-1}\sum_x\rho_N(T_x(A)T_x(\tau_t(B)))\), which equals
\(\omega_N(A\tau_t(B))\); the second line retains the opposite
operator order.  Thus both error bounds tend to zero.

The maps \(t\mapsto A\tau_t(B)\) and
\(t\mapsto\tau_t(B)A\) are norm-continuous.  Their images of a
compact time interval are norm-compact.  Pointwise convergence
\(\omega_{N_j}(C)\to\omega(C)\) is uniform over each such image:
use a finite norm net and the common functional norm one.  Therefore,
along the already chosen subsequence,
\begin{align*}
 F_{N_j}^{A,B}(t)&\longrightarrow\omega(A\tau_t(B)),\\
 F_{N_j}^{A,B}(t+i\beta)&\longrightarrow\omega(\tau_t(B)A)
\end{align*}
locally uniformly for real \(t\).  Lemma~\ref{lem:strip-limit}
now gives a bounded continuous function \(F_{A,B}\) on the closed
strip, analytic in its interior, with precisely these two boundary
values and
\begin{equation}\label{eq:local-kms-bound}
 \sup_{z\in\mathcal S_\beta}|F_{A,B}(z)|\leq\|A\|\|B\|.
\end{equation}
Since the strip limit uses the full state subsequence, the bilinearity of
\eqref{eq:finite-strip-functions} passes to \(F_{A,B}\).

It remains to obtain the asserted condition on the entire algebra.
Let \(A_n,B_n\) be local observables converging in norm to arbitrary
\(A,B\in\mathcal A\).  Bilinearity and
\eqref{eq:local-kms-bound} imply
\[
 \sup_{z\in\mathcal S_\beta}
       |F_{A_n,B_n}(z)-F_{A_m,B_m}(z)|
 \leq \|A_n-A_m\|\|B_n\|
      +\|A_m\|\|B_n-B_m\|.
\]
Hence these functions converge uniformly on the closed strip to a
bounded continuous \(F_{A,B}\), analytic in the interior.  The same
inequality proves independence of the approximating sequences and
the bound \(\sup|F_{A,B}|\leq\|A\|\|B\|\).  Since states have
norm one and \(\tau_t\) is isometric, norm approximation passes
both boundary identities to the limit, uniformly in real \(t\):
\[
 F_{A,B}(t)=\omega(A\tau_t(B)),\qquad
 F_{A,B}(t+i\beta)=\omega(\tau_t(B)A).
\]
This is the full \(\beta\)-KMS condition \eqref{eq:kms-definition}
for the zero-field dynamics \(\tau\) defined in Section~\ref{sec:model}.  The state is translation invariant and satisfies
\(\omega(S_0^z)\geq S/4>0\) by
Proposition~\ref{prop:magnetized-limit}.  Since \(\beta\) was any
fixed finite number at least \(\beta_0(d,S)\), Theorem~\ref{thm:main}
follows.
\end{proof}
